\chapter{Giới thiệu}
\label{chap:introduction}

\section{Bối cảnh và động lực nghiên cứu}

Hiện nay, các mô hình ngôn ngữ lớn (LLM) được sinh viên đại học sử dụng rộng rãi để phác thảo tổng quan tài liệu, tóm tắt bài báo khoa học và định dạng danh mục tài liệu tham khảo. Dù giúp tăng tốc đáng kể quy trình soạn thảo học thuật, các mô hình tạo sinh lại kéo theo một rủi ro nghiêm trọng về độ tin cậy được gọi là \emph{ảo giác trích dẫn}~\cite{Ji2023}, \cite{huang2025hallucination}. Do các mô hình ngôn ngữ tự hồi quy dự đoán chuỗi từ có xác suất cao thay vì truy vấn trực tiếp các danh mục học thuật chính thống, chúng thường xuyên tạo ra những mục tài liệu tham khảo có cú pháp hợp lệ nhưng không vượt qua được các bước kiểm chứng sự thật cơ bản. Các khảo sát thực nghiệm của Walters và Wilder~\cite{walters2023fabrication} cùng Agrawal và cộng sự~\cite{agrawal2024hallucinating} cho thấy một tỷ lệ lớn trích dẫn do LLM sinh ra hoặc không tồn tại trên thực tế, hoặc chứa định danh đối tượng số (DOI) bịa đặt, sai năm xuất bản, hay bị làm sai lệch danh sách tác giả.

Việc kiểm tra thủ công danh mục tài liệu tham khảo không thể đáp ứng quy mô thực tế trong đào tạo đại học. Một báo cáo đồ án tốt nghiệp thông thường trích dẫn từ ba mươi đến tám mươi nguồn tài liệu trải rộng trên tạp chí bình duyệt, kỷ yếu hội nghị, sách chuyên khảo, tài liệu đặc tả kỹ thuật trực tuyến và văn bản quy phạm pháp luật quốc gia. Để rà soát từng mục tham khảo bằng tay, giảng viên phải tra cứu nhiều cơ sở dữ liệu thư mục khác nhau, đối chiếu họ tên tác giả cùng năm xuất bản, kiểm tra xem mã DOI có phân giải đúng về tiêu đề được khai báo hay không, đồng thời đọc phần tóm tắt của nguồn gốc để xác nhận nhận định trong thân bài phản ánh trung thực bài báo gốc. Hơn nữa, lỗi trích dẫn phân tách thành các nhóm khiếm khuyết cấu trúc mang ý nghĩa sư phạm rất khác nhau:
\begin{enumerate}
    \item \textbf{Nguồn tài liệu bịa đặt hoặc lai ghép (\textit{chimeric}):} Bài báo được trích dẫn không tồn tại trong bất kỳ cơ sở dữ liệu học thuật nào, hoặc một mã DOI có thật của bài báo nổi tiếng bị ghép vào một tiêu đề bịa đặt hoặc danh sách tác giả không liên quan.
    \item \textbf{Lỗi siêu dữ liệu hoặc sai lệch nhận định:} Công trình gốc có tồn tại trên thực tế, nhưng bản thảo ghi sai năm xuất bản, bỏ sót đồng tác giả, trích dẫn bài báo đã bị rút lại (\textit{retracted}), hoặc gắn một nhận định trong thân bài mâu thuẫn với phần tóm tắt của tài liệu nguồn.
    \item \textbf{Lỗi liên kết cấu trúc:} Ký hiệu trích dẫn trong thân bài như \textit{(Vaswani et al., 2017)} không có mục tương ứng trong danh mục tài liệu tham khảo, hoặc các mục tham khảo được liệt kê nhằm làm tăng số lượng nguồn nhưng không hề được viện dẫn trong nội dung bài viết.
\end{enumerate}

Hai nguyên tắc kỹ thuật phần mềm cốt lõi định hướng giải pháp xây dựng trong đồ án tốt nghiệp này. Thứ nhất, không một cơ sở dữ liệu thư mục đơn lẻ nào lập chỉ mục được toàn bộ kho tàng xuất bản toàn cầu~\cite{visser2021bibliographic}. Các bài báo hội thảo khoa học máy tính thường xuất hiện trên thư mục khoa học máy tính DBLP (\textit{Digital Bibliography \& Library Project}) hoặc kho lưu trữ của Hiệp hội Ngôn ngữ học Tính toán (ACL Anthology) trước khi được lập chỉ mục trên Crossref, các công trình y sinh tập trung tại PubMed và Europe PubMed Central (PMC), còn báo cáo đại học tại Việt Nam thường xuyên viện dẫn văn bản pháp luật quốc gia (\textit{Luật}, \textit{Nghị định}, \textit{Thông tư}, \textit{Quyết định}). Do đó, kiểm chứng đáng tin cậy đòi hỏi cơ chế truy hồi liên hợp trên nhiều cơ sở dữ liệu bổ trợ lẫn nhau. Thứ hai, một công cụ đánh giá học thuật tự động tuyệt đối không được biến sự cố giao diện lập trình ứng dụng (API) thoáng qua hay kết quả đối sánh giáp ranh thành cáo buộc gian lận sai trái. Tuân theo lý thuyết phân loại có chọn lọc~\cite{chow1970optimum}, \cite{Geifman2017}, hệ thống phải chủ động từ chối phán quyết tường minh khi API bên ngoài giới hạn tần suất hoặc khi bằng chứng đối sánh rơi vào vùng bất định, nhường quyền phán quyết học thuật cuối cùng cho giảng viên.

\section{Phát biểu bài toán và phạm vi nghiên cứu}

Đồ án tốt nghiệp này có tên đề tài là \textit{``CITRA: Nền tảng Nơ-ron -- Ký hiệu phục vụ chấm điểm tự động tính toàn vẹn trích dẫn và phát hiện ảo giác nguồn tài liệu''}. Trong kỹ thuật phần mềm giáo dục, bài toán đánh giá bản thảo tự động được phân tách thành ba hướng tiếp cận độc lập: (i)~\emph{chấm điểm bài luận tổng thể tự động}, vốn ước lượng các thuộc tính chủ quan như bố cục lập luận và văn phong; (ii)~\emph{phát hiện văn bản máy sinh theo phong cách hành văn}, ước lượng khả năng câu chữ bề mặt do trí tuệ nhân tạo (AI) tạo sinh; và (iii)~\emph{kiểm chứng tự động tính toàn vẹn trích dẫn}, tập trung kiểm toán nguồn gốc thư mục khách quan và tính nhất quán cấu trúc. Thay vì theo đuổi chấm luận tổng thể thiếu minh bạch hay nhận diện văn phong mong manh, đồ án giải quyết chiều đánh giá thứ ba có thể kiểm chứng thực nghiệm: \textbf{Chấm điểm tự động tính toàn vẹn trích dẫn (CIS)} và \textbf{Phát hiện ảo giác nguồn tài liệu}.

Phát biểu hình thức, gọi một bản thảo học thuật nộp dưới định dạng Portable Document Format (\texttt{.pdf}), Office Open XML (\texttt{.docx}), BibTeX (\texttt{.bib}) hoặc Research Information Systems (\texttt{.ris}) là bộ ba $D = (T, \mathcal{M}, \mathcal{B})$, trong đó $T$ là văn bản thân bài, $\mathcal{M} = \{m_1, \dots, m_k\}$ là đa tập các ký hiệu trích dẫn nội văn kèm ngữ cảnh câu bao quanh $c(m_i)$, và $\mathcal{B} = \{r_1, \dots, r_N\}$ là danh mục tài liệu tham khảo. Hệ thống phần mềm phải giải quyết bốn tác vụ gắn kết:
\begin{enumerate}
    \item \textbf{Trích xuất đa định dạng và liên kết hai chiều:} Trích xuất $\mathcal{M}$ và $\mathcal{B}$ trên nhiều chuẩn định dạng khác nhau, đồng thời tái tạo ánh xạ hai phía $L: \mathcal{M} \to \mathcal{B} \cup \{\bot\}$ nhận diện tám trạng thái toàn vẹn cấu trúc (bao gồm thiếu tài liệu tham khảo, tài liệu không được trích dẫn và sai lệch năm nội văn).
    \item \textbf{Kiểm chứng liên hợp nguồn tài liệu và nhận định:} Truy vấn các cơ sở dữ liệu cục bộ và trực tuyến, kiểm tra tính khả dụng của định danh cùng trạng thái rút bài, đối sánh ngữ nghĩa nhận định xuyên ngữ với bản tóm tắt chuẩn, và phân loại mỗi mục tham khảo $r_j \in \mathcal{B}$ vào hệ năm nhãn
    \[\begin{gathered}
    \mathcal{Y} = \{\texttt{verified},\; \texttt{metadata\_error},\; \texttt{suspected\_hallucination},\\
    \texttt{unresolved},\; \texttt{resource}\}.
    \end{gathered}\]
    \item \textbf{Chấm điểm toàn vẹn có hiệu chỉnh xác suất và tự động sửa lỗi thư mục:} Gắn cho mỗi phán quyết $y_j \in \mathcal{Y}$ một xác suất độ tin cậy đã hiệu chỉnh $P_{\text{cal},j} \in \{0\} \cup [0.01, 0.99]$ cùng vết luật kiểm toán, tổng hợp bằng chứng toàn tài liệu thành điểm $\mathrm{CIS}(D) \in [0, 100]$ dễ diễn giải, và xuất danh mục tham khảo đã chuẩn hóa từ siêu dữ liệu chính thống.
    \item \textbf{Hậu kiểm có sự tham gia của giảng viên và lưu vết kiểm toán:} Hiển thị mọi kết luận, bản ghi đối chứng và chuỗi suy diễn logic trên không gian web phân quyền an toàn, nơi giảng viên được cấp quyền có thể kiểm tra, ghi đè và truy vết quyết định.
\end{enumerate}

\section{Mục tiêu nghiên cứu}
\label{sec:objectives}

Để giải quyết bài toán trên theo các tiêu chí kỹ thuật và thực nghiệm cụ thể, đo lường được, khả thi, thực tế và có thời hạn (SMART), đồ án xác định sáu mục tiêu gắn kết trực tiếp với năm câu hỏi nghiên cứu (\textbf{RQ1}--\textbf{RQ5}) đánh giá ở Chương~\ref{chap:results}:
\begin{itemize}
    \item \textbf{OBJ-1 (Trích xuất đa định dạng độ phủ cao và liên kết hai chiều --- RQ4):} Xây dựng phân hệ phân tích bản thảo chịu lỗi và liên kết hai chiều đạt độ phủ trích xuất danh mục tham khảo $\ge 95\%$ trên mười hai quy chuẩn trích dẫn học thuật và pháp lý, đồng thời phát hiện đầy đủ tám trạng thái ánh xạ cấu trúc giữa nội văn và danh mục tham khảo.
    \item \textbf{OBJ-2 (Truy hồi nguồn gốc liên hợp nhận biết miền chuyên ngành --- RQ4, RQ5):} Thiết kế kiến trúc truy hồi phân tầng ưu tiên ngoại tuyến đạt độ phủ ngoại tuyến $\ge 99\%$ với độ trễ tra cứu cục bộ dưới 100\,ms trên kho bài báo khoa học máy tính đã lập chỉ mục, kết hợp điều phối truy vấn có giới hạn tốc độ và bộ ngắt mạch (\textit{circuit breaker}) qua các cơ sở dữ liệu học thuật, kỹ thuật và lưu trữ công khai cùng kiểm tra rút bài tự động và xác thực liên kết chống tấn công giả mạo yêu cầu phía máy chủ (SSRF).
    \item \textbf{OBJ-3 (Kiểm chứng lai Nơ-ron -- Ký hiệu giải thích được và hiệu chỉnh xác suất --- RQ1, RQ2, RQ3):} Xây dựng bộ suy diễn lai kết hợp đối sánh ngữ nghĩa song ngữ với các rào chắn nguồn gốc ký hiệu rời rạc và ước lượng xác suất có hiệu chỉnh, hướng tới độ chính xác tầng quyết định $\ge 99\%$ cả trong phân phối và ngoài phân phối (OOD), tỷ lệ kết tội nhầm nghiêm trọng từ hợp lệ sang ảo giác $\le 0.1\%$, và sai số hiệu chỉnh kỳ vọng ($\text{ECE}$) dưới $0.05$ kèm từ chối phán quyết có chọn lọc.
    \item \textbf{OBJ-4 (Chấm điểm toàn vẹn dễ diễn giải và chuẩn hóa thư mục tự động --- RQ3, RQ4):} Định nghĩa thang Điểm toàn vẹn trích dẫn (CIS) giới hạn trên $[0, 100]$ minh bạch, đạt độ ổn định thứ hạng cao (Spearman $\rho \ge 0.98$) dưới nhiễu loạn trọng số, tích hợp xuất và sửa lỗi danh mục tham khảo tự động theo siêu dữ liệu chính thống.
    \item \textbf{OBJ-5 (Nền tảng web phân quyền sẵn sàng vận hành --- RQ5):} Phát triển nền tảng web toàn trình đóng gói bằng container (\textbf{CITRA}) hỗ trợ quy trình Sinh viên, Giảng viên và Quản trị viên, truyền phát tiến độ thời gian thực, ghi đè phán quyết đảm bảo tính giao dịch và phòng thủ chiều sâu về bảo mật, đáp ứng ngân sách độ trễ đầu--cuối ($\le 15\,\text{s}$ cho báo cáo 50 trang) và kích thước gói trình duyệt ($\le 250\,\text{kB}$ cho gói khởi động) trên phần cứng bộ xử lý trung tâm (CPU) phổ thông.
    \item \textbf{OBJ-6 (Đánh giá thực nghiệm quy mô lớn chống rò rỉ dữ liệu --- RQ1--RQ5):} Thực thi quy trình đánh giá có khả năng tái lập và phân tách độc lập theo bài báo gốc, kết hợp bộ chuẩn tầng quyết định 120,000 mẫu (so với sáu mô hình cơ sở \texttt{B0}--\texttt{B5}), bộ kiểm thử rào chắn OOD 2,400 mẫu, đánh giá trực tuyến đa API, kiểm toán bản thảo đa chuẩn thực tế và kiểm thử chịu tải đồng thời với kiểm định ý nghĩa thống kê chuẩn mực.
\end{itemize}

\section{Phạm vi hệ thống và ranh giới vận hành}

\subsection{Các năng lực thuộc phạm vi hệ thống}

Phiên bản \texttt{3.91.3} của \textbf{CITRA} xử lý các bản thảo học thuật tiếng Anh và tiếng Việt ở định dạng \texttt{.pdf}, \texttt{.docx}, \texttt{.bib} và \texttt{.ris}. Tầng trích xuất hỗ trợ mười một chuẩn trích dẫn học thuật quốc tế---Hiệp hội Tâm lý học Hoa Kỳ (APA), Viện Kỹ sư Điện và Điện tử (IEEE), Hiệp hội Máy tính Hoa Kỳ (ACM), Harvard, Vancouver, Nature, Hội đồng Biên tập Khoa học (CSE), Hiệp hội Hóa học Hoa Kỳ (ACS), Hiệp hội Y khoa Hoa Kỳ (AMA), Chicago và Hiệp hội Ngôn ngữ Hiện đại (MLA)---cùng nhóm văn bản quy phạm pháp luật Việt Nam (\textit{Luật}, \textit{Nghị định}, \textit{Thông tư}, \textit{Quyết định}). Phạm vi kiểm chứng nguồn bao phủ bài báo tạp chí, kỷ yếu hội nghị, bản tiền ấn phẩm (\textit{preprint}), sách chuyên khảo, kho lưu trữ truy cập mở, tài liệu đặc tả kỹ thuật (chuẩn Internet IETF RFC, gói thư viện Python, tài nguyên Hugging Face) và tài nguyên web trực tiếp hoặc đã lưu trữ. Toàn bộ thành phần cốt lõi triển khai qua Docker Compose trên máy chủ CPU phổ thông mà không đòi hỏi bộ xử lý đồ họa (GPU) hay giấy phép cơ sở dữ liệu thương mại; các bộ chuyển đổi tìm kiếm web dự phòng trả phí mặc định bị tắt trừ khi đơn vị vận hành chủ động cấu hình khóa truy cập riêng.

\subsection{Các giới hạn nằm ngoài phạm vi}

Năm ranh giới vận hành tường minh xác định phạm vi áp dụng của \textbf{CITRA}:
\begin{enumerate}
    \item \textbf{Không chấm văn phong hay kiểm tra đạo văn kho đóng:} \textbf{CITRA} không chấm chất lượng hành văn, ngữ pháp hay tính nguyên bản của ý tưởng và nội dung bài viết, đồng thời không dò trùng lặp câu chữ trên các kho bài làm sinh viên đóng.
    \item \textbf{Tuân thủ tường phí và kiểm soát truy cập:} Khi một ấn phẩm thương mại nằm sau tường phí của nhà xuất bản và ẩn cả toàn văn lẫn tóm tắt, bước đối sánh ngữ nghĩa nhận định chỉ sử dụng siêu dữ liệu và tóm tắt mở hợp pháp mà không vượt qua cơ chế kiểm soát truy cập.
    \item \textbf{Phạm vi ngôn ngữ và phân bố miền kho cục bộ:} Chuẩn hóa họ tên và văn bản pháp quy chuyên biệt hướng tới bản thảo tiếng Anh và tiếng Việt; tài liệu bằng ngôn ngữ khác sử dụng đối sánh từ vựng và nhúng đa ngữ. Ngoài ra, do kho ngoại tuyến mặc định tập trung vào xử lý ngôn ngữ tự nhiên ($98.42\%$ từ ACL Anthology), các chuyên ngành khác phụ thuộc chủ yếu vào truy vấn trực tuyến đa cơ sở dữ liệu.
    \item \textbf{Vai trò hỗ trợ ra quyết định:} Kết luận của máy là bằng chứng kiểm toán hỗ trợ giảng viên chứ không phải phán quyết kỷ luật tự động; giảng viên luôn giữ toàn quyền thẩm định và ghi đè.
    \item \textbf{Nghiên cứu người dùng dài hạn tại lớp học:} Đánh giá mức giảm tải chấm bài và tác động sư phạm xuyên suốt nhiều học kỳ nằm ngoài khung thời gian đồ án và được định hướng cho giai đoạn triển khai tiếp theo (Chương~\ref{chap:conclusion}).
\end{enumerate}

\section{Đóng góp chính của đồ án}

Bám sát các mục tiêu từ \textbf{OBJ-1} đến \textbf{OBJ-6}, đồ án tốt nghiệp mang lại ba đóng góp về phương pháp luận -- thực nghiệm cùng ba đóng góp về kỹ thuật hệ thống:

\subsection{Đóng góp phương pháp luận và thực nghiệm}
\begin{enumerate}
    \item \textbf{Bộ máy kiểm chứng lai Nơ-ron -- Ký hiệu kèm hiệu chỉnh xác suất Platt:} Kết hợp biểu diễn nhúng Sentence-BERT (SBERT) và bộ mã hóa song ngữ lượng tử hóa theo chuẩn trao đổi mạng nơ-ron mở số nguyên 8-bit (ONNX INT8) với tầng tiền sàng lọc tài nguyên Bậc 0 (Order~0), 16 nhóm luật ký hiệu phân cấp ưu tiên (gồm 68 mã luật \texttt{R-*} thực thi riêng biệt: 65 luật tĩnh và 3 mã điều chỉnh nguồn gốc gắn tiền tố động), mô hình hồi quy logistic đa lớp (MLR) 16 đặc trưng ($C = 0.25$) và hiệu chỉnh xác suất Platt trên không gian logit.
    \item \textbf{Cơ chế từ chối phán quyết có chọn lọc dựa trên nguồn gốc:} Tách biệt rõ ràng giữa bác bỏ có xác nhận từ cơ sở dữ liệu (\texttt{suspected\_hallucination}, đòi hỏi phản hồi HTTP 200 rỗng hợp lệ hoặc định danh bịa đặt/lai ghép) với nghẽn tần suất mạng (\texttt{HTTP 429}), các Cơ quan đăng ký DOI (DOI RA) khu vực không công bố siêu dữ liệu theo Ngôn ngữ định dạng trích dẫn (CSL, \texttt{R-DOI-RA-NO-CSL}) và vùng điểm tương đồng giáp ranh ($[0.40, 0.60)$), chuyển hướng an toàn các trường hợp bất định sang \texttt{unresolved}.
    \item \textbf{Bộ đánh giá thực nghiệm đa tầng có khả năng tái lập:} Thiết lập bộ chuẩn tầng quyết định 120,000 mẫu phân tách độc lập theo bài báo gốc cùng bộ kiểm thử rào chắn OOD 2,400 mẫu kèm quy trình kiểm định bắt cặp McNemar, Holm--Bonferroni và bootstrap (độ chính xác trong phân phối $99.67\%$, $\text{ECE} = 0.0405$, độ chính xác OOD $99.96\%$), bổ sung bởi đánh giá trực tuyến 15 API trên 120 tài liệu tham khảo thô (độ chính xác 5 lớp $89.17\%$; F1 lớp hợp lệ $98.36\%$) và kiểm toán từng dòng trên các bản thảo đa chuẩn thực tế.
\end{enumerate}

\subsection{Đóng góp kỹ thuật hệ thống}
\begin{enumerate}
    \item \textbf{Bộ trích xuất mười hai chuẩn và liên kết hai chiều bảy phương thức:} Kết hợp con trỏ cấu trúc chuẩn mã hóa văn bản học thuật (TEI-XML) từ công cụ GROBID (\textit{GeneRation Of BIbliographic Data}, \texttt{0.9.1-crf}) với cơ chế dự phòng PyMuPDF hai cột, bộ lọc mục lục tất định, bộ kiểm chứng văn bản pháp luật Việt Nam (\texttt{VNLegalVerifier}) và bộ liên kết hai phía phát hiện tám trạng thái ánh xạ cấu trúc.
    \item \textbf{Kiến trúc truy hồi bốn tầng ưu tiên ngoại tuyến chống tấn công SSRF:} Kết hợp bản chụp đánh giá kho tìm kiếm toàn văn cục bộ SQLite FTS5 gồm 129,137 bản ghi (128,115 bài báo duy nhất đủ chuẩn đo kiểm; 82,887 bản tóm tắt) và danh sách 101 công trình kinh điển với định tuyến giới hạn tốc độ có bộ ngắt mạch qua 15+ API công khai cùng bộ kiểm tra HTTP/Wayback ghim địa chỉ IP theo DNS chống SSRF.
    \item \textbf{Nền tảng toàn trình sẵn sàng vận hành (\textbf{CITRA} \texttt{v3.91.3}):} Hiện thực hóa máy chủ FastAPI bất đồng bộ (15 bảng SQLAlchemy ORM, 13 bản di trú lược đồ có kiểm tra mã băm SHA-256 \texttt{0001}--\texttt{0013}, bộ điều phối tiếp nhận \texttt{PipelineAdmissionController} và cơ chế tính lại điểm CIS đảm bảo tính giao dịch) cùng giao diện web React 18.3 / TypeScript 5.9, được xác minh bởi 2,437 ca kiểm thử pytest backend mặc định (2,446 tổng gồm 9 ca chậm/trực tuyến), 495 ca kiểm thử Vitest frontend (63 bộ đặc tả) và 20 bộ kiểm thử đầu--cuối Playwright (178 kịch bản: 174 E2E Chromium + 4 kiểm khói sản phẩm trực tiếp; tổng cộng 3,110 kiểm tra mặc định, 3,119 tổng).
\end{enumerate}

\section{Cấu trúc báo cáo đồ án}

Báo cáo đồ án tốt nghiệp này được tổ chức thành sáu chương và năm phụ lục (Phụ lục~\ref{app:config}--\ref{app:log}), tóm tắt trong Bảng~\ref{tab:outline}.

\begin{table}[!htbp]
\centering
\caption[Cấu trúc báo cáo đồ án tốt nghiệp]{Cấu trúc báo cáo đồ án tốt nghiệp theo từng chương và phụ lục.}\label{tab:outline}
\begin{tabular}{|L{0.24}|L{0.76}|}
\hline
\textbf{Phần} & \textbf{Tóm tắt nội dung} \\
\hline
\textbf{Chương 1} & Bối cảnh, động lực nghiên cứu, phát biểu bài toán hình thức, mục tiêu SMART (\textbf{OBJ-1}--\textbf{OBJ-6}), phạm vi vận hành và đóng góp chính. \\
\hline
\textbf{Chương 2} & Tổng quan tài liệu (gồm công trình 2024--2025), đối sánh khách quan với công cụ hiện có, cơ sở dữ liệu liên hợp và nền tảng toán học. \\
\hline
\textbf{Chương 3} & Yêu cầu chức năng IEEE 29148 và phi chức năng ISO/IEC 25010, sơ đồ UML 2.5 (ca sử dụng, tuần tự, triển khai), ERD \textbf{15 bảng}, API và bảo mật. \\
\hline
\textbf{Chương 4} & Hiện thực hóa \textbf{tám tầng phần mềm}, mã giả thuật toán, bộ suy diễn ký hiệu \textbf{68 luật}, mô hình 16D + Platt, công thức CIS và giao diện React. \\
\hline
\textbf{Chương 5} & Truy vết yêu cầu, kiểm thử tự động, thực nghiệm trả lời \textbf{RQ1}--\textbf{RQ5} (tầng quyết định, API trực tuyến, PDF thực, chịu tải) và mối đe dọa độ giá trị. \\
\hline
\textbf{Chương 6} & Đối chiếu mức độ đạt mục tiêu, tổng hợp các hạn chế hiện hữu và định hướng nghiên cứu, phát triển trong tương lai. \\
\hline
\textbf{Phụ lục~\ref{app:config}--\ref{app:log}} & Cấu hình môi trường (\ref{app:config}) và Docker Compose (\ref{app:docker}), bảng phân loại năm nhãn (\ref{app:taxonomy}), sơ đồ gói mã nguồn (\ref{app:packages}) và nhật ký thực hiện (\ref{app:log}). \\
\hline
\end{tabular}
\end{table}
